% ==========================================================================
% ORAL — PARTIE ENRIQUE (slides 1 a 10)
% Script precis : tout ce qu'il faut dire, mot pour mot.
% Compilable avec pdflatex.
% ==========================================================================
\documentclass[11pt,a4paper]{article}

\usepackage[utf8]{inputenc}
\usepackage[T1]{fontenc}
\usepackage[french]{babel}
\usepackage{lmodern}
\usepackage[a4paper,margin=2.2cm]{geometry}

\usepackage{amsmath,amssymb,bm}
\usepackage{xcolor}
\usepackage{hyperref}
\usepackage{tcolorbox}
\usepackage{enumitem}
\usepackage{booktabs}
\usepackage{array}
\usepackage{fancyhdr}
\usepackage{listings}

\hypersetup{colorlinks=true, linkcolor=blue!60!black, urlcolor=blue!60!black}

% Style des blocs de code Python
\definecolor{codebg}{rgb}{0.97,0.97,0.97}
\definecolor{codekw}{rgb}{0.13,0.29,0.53}
\definecolor{codestr}{rgb}{0.50,0.20,0.00}
\definecolor{codecom}{rgb}{0.25,0.50,0.25}
\lstdefinestyle{py}{
    language=Python,
    backgroundcolor=\color{codebg},
    basicstyle=\ttfamily\footnotesize,
    keywordstyle=\color{codekw}\bfseries,
    stringstyle=\color{codestr},
    commentstyle=\color{codecom}\itshape,
    showstringspaces=false,
    breaklines=true,
    frame=single,
    rulecolor=\color{gray!40},
    numbers=left,
    numberstyle=\tiny\color{gray},
    numbersep=6pt,
    captionpos=b,
    tabsize=4,
    columns=fullflexible
}
\lstset{style=py}

\setlength{\parindent}{0pt}
\setlength{\parskip}{6pt plus 1pt}

\pagestyle{fancy}
\fancyhf{}
\fancyhead[L]{Oral — Partie Enrique}
\fancyhead[R]{Battle Météorage}
\fancyfoot[C]{\thepage}

\newtcolorbox{dire}{colback=blue!4, colframe=blue!50!black, title=\textbf{À dire (mot pour mot)}}
\newtcolorbox{compr}{colback=green!4, colframe=green!50!black, title=\textbf{Ce qu'il faut comprendre}}
\newtcolorbox{piege}{colback=orange!4, colframe=orange!60!black, title=\textbf{Piège / question possible du jury}}
\newtcolorbox{transit}{colback=gray!5, colframe=gray!50!black, title=\textbf{Transition}}

\title{\textbf{Oral — Partie Enrique} \\[0.3em]
       \large Script précis pour les slides 1 à 10 \\
       \small Battle Météorage --- Weibull AFT + protocole Data Battle}
\author{}
\date{}

% ==========================================================================
\begin{document}
\maketitle

\begin{tcolorbox}[colback=yellow!8, colframe=yellow!50!black, title=\textbf{Mode d'emploi}]
\textbf{Tu présentes les slides 1 à 10.} Pour chaque slide :
\begin{itemize}[leftmargin=*, topsep=2pt]
    \item \textbf{À dire} = phrase à dire telle quelle (en bleu)
    \item \textbf{Ce qu'il faut comprendre} = ce que le concept veut dire en français simple (en vert)
    \item \textbf{Piège} = question pénible que le jury peut poser, et la réponse
    \item \textbf{Transition} = phrase pour enchaîner vers le slide suivant (en gris)
\end{itemize}

Lis les blocs bleus en gardant les yeux sur le public. Les autres sont pour toi en cas de question.

\textbf{Important} : si tu galères sur un slide, dis juste les blocs bleus, ils suffisent.
Les autres sont du bonus pour répondre aux questions.
\end{tcolorbox}

\tableofcontents
\newpage

% ==========================================================================
\section{Slide 0 — Titre : Battle Météorage}

\begin{dire}
\og Bonjour. Notre sujet : la \textbf{Battle Météorage}.

Peut-on raccourcir la règle des 30 minutes d'alerte orage actuellement appliquée par Météorage,
sans augmenter le risque pour le personnel au sol des aéroports ?

Notre réponse : un modèle \textbf{Weibull AFT calibré}, validé par une vérification empirique
indépendante, qui adapte le temps d'attente au contexte de chaque éclair.

Pendant cette présentation je vais d'abord poser le problème, présenter nos données et notre
modèle, puis mon binôme prendra le relais pour les résultats et la démo. \fg{}
\end{dire}

\begin{compr}
\textbf{Météorage} = entreprise française qui détecte les éclairs en temps réel par triangulation
d'ondes radio. Elle équipe la majorité des aéroports européens en service d'alerte.

\textbf{Battle Météorage} = défi 2026 de Pau IA. L'objectif : améliorer la règle actuelle
\og 30 min après le dernier éclair \fg{} en utilisant l'IA.

\textbf{Weibull AFT calibré} = un modèle de survie probabiliste (on prédit \emph{combien de
temps} jusqu'au prochain événement) recalibré empiriquement sur les données pour que les
probabilités annoncées soient fiables.
\end{compr}

\begin{transit}
\og Voyons d'abord le problème. \fg{}
\end{transit}

% ==========================================================================
\section{Slide 1 — Le problème : la règle 30 min est aveugle au contexte}

\begin{dire}
\og Aujourd'hui, dès qu'un éclair \textbf{nuage-sol} (qu'on appelle CG, pour
\emph{cloud-to-ground}) tombe dans un rayon de 20 km autour d'un aéroport, Météorage
\textbf{déclenche une alerte}.

Pendant cette alerte, l'aéroport interrompt ses opérations au sol : atterrissages, ravitaillements,
chargements. C'est essentiel pour la sécurité du personnel et des aéronefs au sol.

L'alerte est \textbf{levée 30 minutes après le dernier CG observé}. Chaque nouveau CG remet
le compteur à zéro.

Le problème : cette règle de 30 minutes est \textbf{aveugle au contexte}. Tous les orages
reçoivent le même délai forfaitaire. Or :
\begin{itemize}
    \item Trop court $\Rightarrow$ risque de foudroiement du personnel ou d'un aéronef
    $\Rightarrow$ inacceptable.
    \item Trop long $\Rightarrow$ retards de vols, immobilisation au sol $\Rightarrow$ coûteux
    pour l'aéroport.
\end{itemize}

Notre question : tous les orages ont-ils \emph{besoin} des 30 min de marge ?
\fg{}
\end{dire}

\begin{compr}
\textbf{CG vs IC} :
\begin{itemize}
    \item \textbf{CG} (Cloud-to-Ground) = éclair qui touche le sol. Peut foudroyer un
    technicien, abîmer un avion. \emph{Seuls les CG comptent pour Météorage et ferment
    l'alerte}.
    \item \textbf{IC} (Intra-Cloud) = éclair qui reste dans le nuage. Visible mais non
    dangereux au sol. \emph{Pas pris en compte pour fermer l'alerte}.
\end{itemize}

\textbf{Pourquoi 30 min ?} C'est le délai standardisé par Météorage. On verra plus tard
(slide \og vérification empirique \fg{}) que ce 30 min correspond statistiquement à
\textbf{1 \% de risque résiduel} sur les données train — donc la règle n'est pas
arbitraire, elle est calibrée pour ce niveau de sécurité.

\textbf{Le compromis} : 30 min est conservateur. Pour certains orages, on pourrait lever plus
tôt sans danger. Pour d'autres (orages dense, encore actifs), 30 min ne suffirait peut-être pas.
\end{compr}

\begin{piege}
\textbf{Q : Pourquoi pas juste réduire 30 à 20 ou 15 min pour tout le monde ?}

R : Parce que c'est aveugle au contexte. Un orage de fin de soirée qui s'éteint pourrait
être levé à 10 min. Un orage dense en plein milieu d'un système orageux pourrait nécessiter
40 min. Le forfait global ne capture pas ces différences. Notre modèle, lui, adapte le
délai \emph{par éclair}.
\end{piege}

\begin{transit}
\og Notre approche pour résoudre ce problème : un modèle qui adapte le temps d'attente. \fg{}
\end{transit}

% ==========================================================================
\section{Slide 2 — Notre approche : prédire un temps d'attente adapté}

\begin{dire}
\og Notre approche en trois étapes.

\textbf{1. On observe} le contexte de chaque éclair : son rythme passé (les gaps entre éclairs
précédents), sa distance à l'aéroport, son intensité, la saison.

\textbf{2. Un modèle Weibull AFT} estime la distribution du temps jusqu'au prochain CG.

\textbf{3. On décide} combien de minutes attendre selon un niveau de risque cible.

Le client choisit son niveau de risque acceptable, et on lui donne le temps d'attente
correspondant :
\begin{itemize}
    \item Risque 10 \% (agressif) : on lève vite, on accepte 1 chance sur 10 de manquer un éclair.
    \item Risque 5 \% (compromis) : 1 chance sur 20.
    \item Risque 1 \% (sécurité max) : 1 chance sur 100.
\end{itemize}

Ce paramètre s'appelle $q$ dans notre modèle. \fg{}
\end{dire}

\begin{compr}
\textbf{Pourquoi un modèle de survie ?} Parce que la question \og combien de temps jusqu'au
prochain événement ? \fg{} est exactement le terrain de jeu de l'analyse de survie (née en
médecine pour estimer les durées avant rechute / décès).

\textbf{Weibull AFT} = un type spécifique de modèle de survie. Le \og W \fg{} pour Weibull
(la loi de probabilité utilisée), \og AFT \fg{} pour \textit{Accelerated Failure Time}
(comment les variables explicatives modifient le temps caractéristique). On détaille slide 4.

\textbf{$q$ = niveau de couverture} : si $q = 0{,}99$, on prédit un temps $T_q$ tel que
\emph{avec 99 \% de probabilité}, aucun éclair n'arrive avant ce $T_q$. Donc 1 \% de risque
résiduel.
\end{compr}

\begin{piege}
\textbf{Q : Pourquoi pas un réseau de neurones ou un gradient boosting ?}

R : Trois raisons. (1) La \textbf{censure native} : le dernier CG d'une alerte n'a pas de
prochain CG observable (on sait juste qu'aucun n'est arrivé pendant 30 min). Les GBM/NN
classiques ne savent pas gérer ça nativement, Weibull AFT oui. (2) Les \textbf{quantiles
analytiques} : on a une formule fermée pour $T_q$. (3) L'\textbf{interprétabilité} : on lit
directement l'effet de chaque feature via son coefficient $\beta$.
\end{piege}

\begin{transit}
\og Présentons les données qu'on a utilisées. \fg{}
\end{transit}

% ==========================================================================
\section{Slide 3 — Les données : 5 aéroports, 54k CG, 2400 alertes}

\begin{dire}
\og Les données nous viennent directement de Météorage : un fichier
\texttt{segment\_alerts\_all\_airports\_train.csv} contenant 507 071 éclairs entre 2016 et
2022, sur 5 aéroports (Ajaccio, Bastia, Biarritz, Nantes, Pise).

Chaque ligne représente un éclair détecté dans un rayon de 30 km d'un aéroport, avec :
sa date, sa position (longitude, latitude), son amplitude (intensité du courant), sa distance
à l'aéroport, son type (CG ou IC).

Après filtrage — on ne garde que les CG dans la zone d'alerte 20 km, et on exclut Pise 2016
qui utilisait un système d'enregistrement différent — il nous reste \textbf{54 470 CG}
répartis sur \textbf{2 400 alertes}.

Le point clé pour le modèle : la \textbf{définition d'une alerte}. Un CG dans 20 km démarre
une alerte. L'alerte est une série de CG espacés de moins de 30 min. Elle se ferme 30 min
après le dernier CG. Un nouveau CG après ces 30 min démarre une \emph{nouvelle} alerte.

Implication essentielle : dans une même alerte, le gap entre deux CG consécutifs est
nécessairement inférieur à 30 min. \textbf{On ne peut donc jamais observer un gap intra-alerte
supérieur à 30 min}. C'est ce qu'on appelle en analyse de survie de la \textbf{censure à droite},
et notre modèle Weibull AFT la gère nativement. \fg{}
\end{dire}

\begin{compr}
\textbf{Censure à droite} : on connaît une \emph{borne inférieure} de la durée, mais pas la
vraie valeur.

Exemple : un dernier CG d'une alerte. On sait qu'aucun CG ne suit dans les 30 min suivantes
(sinon il serait dans la même alerte). Mais on ne sait pas \emph{combien de temps après ça}
le prochain CG arrive (45 min ? 2 heures ? 1 semaine ?). La durée jusqu'au prochain CG est
\og censurée à droite à 30 min \fg{}.

\textbf{Pourquoi on exclut Pise 2016 ?} Le capteur de Pise enregistrait différemment les IC
en 2016 vs après. Pour avoir des données homogènes, on l'écarte. Décision prise dans le
notebook d'exploration.

\textbf{Pourquoi 20 km et pas 30 ?} 20 km = rayon officiel de la zone d'alerte Météorage.
Au-delà, les éclairs ne déclenchent pas d'alerte. Le fichier source contient quand même les
éclairs jusqu'à 30 km pour des analyses étendues.
\end{compr}

\begin{piege}
\textbf{Q : Et les données 2023-25 ?}

R : Elles servent à l'évaluation finale (eval set). Le modèle est entraîné uniquement sur
2016-2022 et évalué sur 2023-25, sans data leakage. Mon binôme détaillera les résultats sur
ces données.

\textbf{Q : Pourquoi seulement 5 aéroports ?}

R : Ce sont ceux fournis par Météorage dans le dataset officiel du Data Battle. En production,
on pourrait étendre à d'autres aéroports avec re-entraînement sur des données locales.
\end{piege}

\begin{transit}
\og Maintenant la théorie du modèle Weibull AFT. \fg{}
\end{transit}

% ==========================================================================
\section{Slide 4 — Le modèle Weibull AFT : les formules}

\begin{dire}
\og La loi de Weibull est paramétrée par deux nombres : $\lambda$ (échelle) et $k$ (forme).

La \textbf{fonction de survie} — la probabilité d'attendre \emph{plus} de $t$ minutes — vaut :
\[
S(t \mid \lambda, k) = \exp\!\left(-(t/\lambda)^k\right)
\]

Et le \textbf{hazard} — la probabilité instantanée qu'un éclair arrive maintenant — vaut :
\[
h(t \mid \lambda, k) = \frac{k}{\lambda}\,(t/\lambda)^{k-1}
\]

\textbf{Sur nos données, on trouve $k = 0{,}697$, donc $k < 1$}. Ça veut dire que le hazard
décroît avec le temps : plus on attend longtemps depuis le dernier CG, moins il est probable
qu'un nouveau CG arrive. C'est cohérent avec l'intuition physique d'un orage qui s'éteint
progressivement.

Maintenant le \textbf{modèle AFT} : pour personnaliser la prédiction par éclair, on relie le
paramètre d'échelle $\lambda$ aux 17 variables contextuelles via :
\[
\log \lambda(X) = \beta_0 + \beta_1 X_1 + \dots + \beta_{17} X_{17}
\]

Chaque coefficient $\beta_j > 0$ allonge le temps caractéristique. $\beta_j < 0$ le raccourcit.

Et finalement, le \textbf{temps recommandé} $T_q$ pour un niveau $q$ a une formule fermée :
\[
T_q(X) = \lambda(X) \cdot (-\ln q)^{1/k}
\]
\fg{}
\end{dire}

\begin{compr}
\textbf{$\lambda$ (échelle, en minutes)} : le temps caractéristique. Si $\lambda = 2$, on
attend en moyenne autour de 2 minutes le prochain CG.

\textbf{$k$ (forme, sans unité)} : gouverne l'évolution du hazard dans le temps.
\begin{itemize}
    \item $k < 1$ (notre cas) : hazard décroissant. Plus on attend, moins on a de chances qu'un
    nouvel éclair arrive. \emph{Orage qui s'éteint.}
    \item $k = 1$ : hazard constant. Loi exponentielle, sans mémoire (peu réaliste ici).
    \item $k > 1$ : hazard croissant. Plus on attend, plus la prochaine occurrence est probable.
    \emph{Vieillissement} (typique en fiabilité industrielle).
\end{itemize}

\textbf{Hazard vs survie} : c'est juste deux façons de voir la même chose. $h(t)$ dit \og quel est
le taux de défaillance maintenant ? \fg{}. $S(t)$ dit \og quelle est la proba d'avoir survécu
jusqu'à $t$ ? \fg{}. Reliés par $h = -d\log S / dt$.

\textbf{$T_q$ — la formule magique} : si je veux que $S(t) = q$ (par exemple 99 \% de chances
de survie), je cherche le $t$ tel que $\exp(-(t/\lambda)^k) = q$, ce qui donne
$t = \lambda \cdot (-\ln q)^{1/k}$. C'est le temps recommandé.
\end{compr}

\begin{piege}
\textbf{Q : Pourquoi avoir choisi Weibull et pas une autre loi (Gamma, log-normale) ?}

R : Trois raisons : (1) Weibull est le standard en analyse de survie pour les inter-arrivées,
référence dans les manuels de fiabilité industrielle. (2) lifelines le supporte nativement
avec gestion de la censure. (3) Les quantiles ont une formule analytique simple. On a quand
même testé d'autres lois (cf.\ section calibration plus tard) : Weibull avec calibration
empirique reste le meilleur compromis.

\textbf{Q : Comment lifelines estime $k$ et les $\beta$ ?}

R : Par maximum de vraisemblance — c'est exactement le sujet du slide suivant.
\end{piege}

\begin{transit}
\og Comment on ajuste ces paramètres ? Par maximum de vraisemblance. \fg{}
\end{transit}

% ==========================================================================
\section{Slide 5 — Apprentissage par maximum de vraisemblance avec censure}

\begin{dire}
\og On ajuste $\beta$ et $k$ pour maximiser la \textbf{probabilité des données observées}
selon le modèle. C'est le maximum de vraisemblance, MLE en anglais.

Pour chaque éclair, on distingue \textbf{deux cas} :

\textbf{Cas 1 — éclair non-dernier} (variable d'événement $\delta_i = 1$). On a observé exactement
le temps $t_i$ jusqu'au prochain CG. Le modèle doit donner une \emph{forte densité} en ce point :
\[
\mathcal{L}_i = f(t_i \mid X_i)
\]

\textbf{Cas 2 — dernier éclair d'une alerte} (variable $\delta_i = 0$, c'est la censure).
On sait juste qu'aucun éclair n'est arrivé pendant 30 minutes. Le modèle doit donner une
\emph{forte survie à 30 min} :
\[
\mathcal{L}_i = S(30 \mid X_i)
\]

Et la log-vraisemblance totale qu'on maximise :
\[
\ell(\beta, k) = \sum_{i=1}^{n} \Big[ \delta_i \log f(t_i \mid X_i) + (1 - \delta_i) \log S(30 \mid X_i) \Big]
\]

C'est une fonction concave dans nos paramètres avec régularisation L2, on l'optimise par
Newton-Raphson via la bibliothèque lifelines en 3 lignes de Python. \fg{}
\end{dire}

\begin{compr}
\textbf{MLE — maximum de vraisemblance} : on cherche les paramètres qui rendent les données
observées \emph{les plus probables} selon le modèle. C'est une procédure classique en
statistique paramétrique.

\textbf{Gérer la censure} : la subtilité ici. Un éclair non-dernier contribue avec la
\emph{densité} $f(t_i)$ — \og je vois précisément le prochain CG arriver à $t_i$, le modèle
doit accorder une forte chance à cette valeur exacte \fg{}. Un éclair dernier contribue avec
la \emph{survie} $S(30)$ — \og je sais juste qu'aucun n'est arrivé pendant 30 min, le modèle
doit accorder une forte chance à cette absence \fg{}.

\textbf{Le code} : 3 lignes via lifelines. Sous le capot, l'optimisation utilise
Newton-Raphson avec une régularisation ridge \texttt{penalizer=0.05} qui stabilise le fit
en présence de features corrélées.
\end{compr}

\begin{piege}
\textbf{Q : Sans gérer la censure, qu'est-ce qui se passerait ?}

R : Le modèle serait \emph{biaisé vers le bas} — il sous-estimerait les longs silences
puisqu'il ne les voit jamais directement. La censure permet d'intégrer cette information \og
absence d'observation \fg{} dans la vraisemblance. C'est la grande force de l'analyse de survie.

\textbf{Q : Pourquoi penalizer = 0.05 et pas 0 ?}

R : Régularisation L2 ridge ajoutée à la log-vraisemblance ($\lambda \sum \beta_j^2$). À 0,
le modèle est plus instable numériquement avec 17 features + 4 dummies. À 0.1, il est trop
régularisé et perd un peu de pouvoir prédictif. 0.05 est le bon compromis trouvé par CV rapide.
\end{piege}

\begin{transit}
\og Voyons les 17 variables qu'on a construites pour caractériser chaque éclair. \fg{}
\end{transit}

% ==========================================================================
\section{Slide 6 — Feature engineering : 17 variables contextuelles}

\begin{dire}
\og On construit 17 variables qui résument le \textbf{passé de l'orage} au moment de chaque
éclair. Aucune ne voit le futur — pas de fuite de données.

On les regroupe en 4 catégories :

\textbf{Rythme (6 variables)} : \texttt{cg\_rank}, \texttt{time\_since\_start},
\texttt{prev\_gap}, \texttt{rolling\_gap\_3}, \texttt{rolling\_gap\_5}, \texttt{trend\_gap}.
Elles encodent comment l'orage s'est comporté jusqu'à présent. C'est notre signal le plus fort.

\textbf{Distance (3 variables)} : \texttt{dist\_current}, \texttt{dist\_cum\_mean},
\texttt{dist\_diff}. Pour savoir si l'orage s'éloigne ou s'approche.

\textbf{Intensité (4 variables)} : \texttt{amp\_abs}, \texttt{amp\_cum\_mean},
\texttt{is\_negative}, \texttt{pct\_neg\_cum}. Polarité et puissance des éclairs.

\textbf{Saisonnalité (4 variables)} : \texttt{month\_sin}, \texttt{month\_cos},
\texttt{hour\_sin}, \texttt{hour\_cos}. Encodage cyclique pour capturer les variations
saisonnières et journalières.

La variable la plus importante : \texttt{rolling\_gap\_5} — la moyenne mobile des intervalles
entre les 5 derniers CG. Petit ($<1$ min) = orage dense, encore actif, $T_q$ court.
Grand ($>4$ min) = orage qui se calme, $T_q$ long.

Toutes les variables continues sont standardisées (moyenne 0, écart-type 1) avant le fit. \fg{}
\end{dire}

\begin{compr}
\textbf{Encodage cyclique} : pour les variables temporelles (mois, heure), on utilise $\sin$ et
$\cos$ au lieu d'un encodage ordinaire. Pourquoi ? Pour éviter la fausse discontinuité entre
décembre (mois 12) et janvier (mois 1) — qui sont en réalité voisins dans le cycle annuel.

$\texttt{month\_sin} = \sin(2\pi \cdot \texttt{mois}/12)$,
$\texttt{month\_cos} = \cos(2\pi \cdot \texttt{mois}/12)$.

\textbf{Pas de leak} : toutes les variables utilisent uniquement les \emph{CG déjà arrivés}.
\texttt{rolling\_gap\_5} prend la moyenne des 5 \emph{derniers} gaps, jamais des futurs. C'est
garanti par la construction avec \texttt{groupby + transform} en pandas.

\textbf{Standardisation} : indispensable pour qu'aucune feature ne domine artificiellement (ex.
\texttt{time\_since\_start} en minutes peut atteindre 100+ alors que \texttt{is\_negative}
vaut 0 ou 1). Le scaler est fitté uniquement sur le train (jamais sur le test) pour éviter
toute fuite.
\end{compr}

\begin{piege}
\textbf{Q : Pourquoi pas plus de variables (gradient des amplitudes, fréquence locale, etc.) ?}

R : On a expérimenté. Au-delà de ces 17, le gain de C-index est marginal et le risque
d'overfitting augmente. On a privilégié la parcimonie et l'interprétabilité.

\textbf{Q : Comment on gère le tout premier CG d'une alerte (qui n'a pas de \texttt{prev\_gap}) ?}

R : On met \texttt{prev\_gap = 0} et \texttt{cg\_rank = 1}. Ces cas sont rares (1/N par alerte)
et leur contribution au fit est faible. Une alternative serait d'utiliser le temps depuis la
fin de l'alerte précédente, mais on ne l'a pas exploité ici.
\end{piege}

\begin{transit}
\og Mais est-ce que ce modèle apporte vraiment quelque chose vs une simple moyenne ? Slide 7. \fg{}
\end{transit}

% ==========================================================================
\section{Slide 7 — Pourquoi pas juste une moyenne par aéroport ?}

\begin{dire}
\og C'est la question naturelle : pourquoi ne pas juste prendre la médiane des gaps observés
par aéroport, et l'appliquer comme temps d'attente ?

On a comparé 5 approches, du plus simple au plus sophistiqué, via le \textbf{C-index} sur le
jeu de test.

Le C-index, ou \emph{Concordance Index}, mesure la capacité d'un modèle à \textbf{ordonner}
deux observations selon leur vraie durée. On tire deux observations au hasard, si le modèle
prédit que la première durera plus longtemps et c'est effectivement le cas, on compte 1.
Sinon 0. C-index = moyenne sur toutes les paires.
\begin{itemize}
    \item C-index = 0,50 : aléatoire, pile/face
    \item C-index = 1,00 : ordre parfait
    \item C-index $\geq$ 0,70 : bon modèle
\end{itemize}

Les résultats :
\begin{itemize}
    \item Constant (médiane globale) : 0,50 — aléatoire
    \item Médiane par aéroport : 0,55 — distingue les aéroports mais rien à l'intérieur
    \item P95 empirique par aéroport : 0,54 — biaisé par la censure
    \item Weibull AFT, dummies aéroport seuls (sans features) : 0,55
    \item \textbf{Weibull AFT + 17 features contextuelles : 0,74}
\end{itemize}

\textbf{+20 points de C-index} grâce aux features contextuelles. C'est ÇA le bénéfice du modèle
par rapport à une moyenne. \fg{}
\end{dire}

\begin{compr}
\textbf{C-index = vitesse d'ordonnancement} : c'est l'AUC de la prédiction de survie. Très
courant en survival analysis.

Pourquoi les baselines simples plafonnent à 0,55 ? Parce qu'elles ne discriminent que par
\emph{aéroport}. Tous les éclairs d'Ajaccio reçoivent la même prédiction, idem pour Bastia, etc.
Au sein d'un même aéroport, elles ne savent pas distinguer un orage actif d'un orage qui se
calme.

Notre modèle, par contre, prédit \emph{par éclair} selon le contexte : si \texttt{rolling\_gap\_5}
est petit, T\_q sera court. Si grand, T\_q sera long. D'où le saut de C-index de 0,55 à 0,74.

\textbf{Pourquoi pas 0,99 ?} Parce qu'il reste de la variabilité intrinsèque aux orages (deux
contextes identiques peuvent donner deux gaps différents par hasard). 0,74 est un très bon score
sur des données aussi bruitées.
\end{compr}

\begin{piege}
\textbf{Q : Et un C-index sur le train, c'est combien ?}

R : 0,73 sur le train (cf. notebook). Train et test sont très proches, donc pas d'overfitting.
La régularisation L2 fait son travail.

\textbf{Q : Pourquoi P95 par aéroport ne marche pas mieux ?}

R : Parce que la censure à 30 min biaise les quantiles empiriques. On ne peut jamais observer
de gaps > 30 min intra-alerte. Donc le P95 empirique \emph{sous-estime} la vraie queue. C'est
exactement ce que Weibull AFT corrige nativement via la log-vraisemblance avec censure.
\end{piege}

\begin{transit}
\og Et pour illustrer concrètement, voici deux orages réels au même aéroport. \fg{}
\end{transit}

% ==========================================================================
\section{Slide 8 — Exemple concret : 2 orages au même aéroport (Ajaccio)}

\begin{dire}
\og Pour rendre ça concret, voici deux orages réels enregistrés à Ajaccio. Même aéroport,
même saison — mais contextes très différents au moment du dernier CG observé.

\textbf{Orage A — encore actif.} Le dernier CG observé est dans un orage dense. Son contexte :
\begin{itemize}
    \item \texttt{rolling\_gap\_5 = 2,07 min} (gaps moyens petits, orage très actif)
    \item Distance à l'aéroport = 18,9 km (l'orage est encore loin mais en zone d'alerte)
    \item \texttt{cg\_rank = 49} (49 CG observés dans cette alerte = orage long)
\end{itemize}
Notre Weibull AFT prédit $T_{99} = 18{,}6$ min.

\textbf{Orage B — qui se calme.} Le dernier CG est dans un orage qui s'éteint. Contexte :
\begin{itemize}
    \item \texttt{rolling\_gap\_5 = 5,53 min} (gaps moyens grands, orage lent)
    \item Distance = 10,0 km (proche)
    \item \texttt{cg\_rank = 12} (orage plus court)
\end{itemize}
Notre Weibull AFT prédit $T_{99} = 99{,}7$ min.

\textbf{Comparaison} :
\begin{itemize}
    \item Le modèle Weibull adapte parfaitement : 18,6 min pour A (orage actif, faut être
    prudent), 99,7 min pour B (orage calme, on peut attendre longtemps si on veut 1 \% de risque).
    \item \textbf{Une moyenne par aéroport donnerait 0,48 min pour les DEUX orages} — incapable
    de distinguer.
\end{itemize}

C'est ÇA la valeur ajoutée du contexte. \fg{}
\end{dire}

\begin{compr}
\textbf{Ce que cet exemple montre} :
\begin{itemize}
    \item \textbf{Orage A} : actif, dense, beaucoup d'éclairs récents → le modèle dit \og continue
    d'attendre, le prochain CG est probablement proche \fg{}. $T_{99} = 18{,}6$ min est court
    \emph{pour la couverture 99 \%}.
    \item \textbf{Orage B} : calme, peu d'éclairs récents → le modèle dit \og ça s'éteint, mais
    pour avoir 99 \% de chance de \emph{vraiment} ne plus avoir d'éclair, il faut attendre presque
    100 min \fg{}.
\end{itemize}

\textbf{Pourquoi B est LONG ?} Parce que pour un orage calme, la queue de la distribution est très
étalée — il y a une petite chance qu'un éclair tardif arrive. Pour couvrir 99 \% de cas, il faut
aller loin dans la queue → $T_{99}$ grand.

Tandis qu'un orage actif a une queue resserrée (les éclairs viennent vite ou pas), donc $T_{99}$
plus court.

\textbf{Pourquoi la médiane par aéroport est si basse (0,48 min) ?} Parce que la médiane prend
le milieu de la distribution observée. Or sur Ajaccio, la \emph{plupart} des gaps intra-orage
sont très courts (moins d'une minute). La médiane reflète ce comportement moyen, mais ne
capture pas du tout les longs silences.
\end{compr}

\begin{piege}
\textbf{Q : Est-ce que ces deux orages sont des cas pathologiques ?}

R : Non, ils sont représentatifs des deux régimes typiques : orage actif (rythme rapide) vs
orage qui s'éteint (rythme lent). On les a sélectionnés justement parce qu'ils sont
représentatifs, pas extrêmes. On a vérifié sur l'ensemble du test que la distribution des
$T_q$ prédits est très étalée — preuve que le modèle adapte vraiment au contexte.

\textbf{Q : 99,7 min de prédiction, mais on cape à 30 min en pratique non ?}

R : Exact. En opérationnel, on cape $T_q$ à 30 min — c'est notre filet de sécurité pour ne
jamais être moins prudent que la baseline Météorage. Ici les 99,7 min montrent la prédiction
\emph{brute} du modèle, avant cap.
\end{piege}

\begin{transit}
\og Et pour quantifier précisément quelles variables contribuent le plus, slide 9. \fg{}
\end{transit}

% ==========================================================================
\section{Slide 9 — Le rythme concentre 57 \% du pouvoir prédictif}

\begin{dire}
\og On peut quantifier l'importance de chaque catégorie de variables en sommant la valeur
absolue des coefficients $\beta$ standardisés.

Résultats sur notre modèle :
\begin{itemize}
    \item \textbf{Rythme} : \textbf{57 \%} du pouvoir prédictif total. Ce sont
    \texttt{rolling\_gap\_5} et \texttt{rolling\_gap\_3} principalement, plus \texttt{cg\_rank}
    et \texttt{prev\_gap}.
    \item \textbf{Saisonnalité} : 15 \%. \texttt{month\_sin}, \texttt{month\_cos},
    \texttt{hour\_sin}, \texttt{hour\_cos}.
    \item \textbf{Aéroport} (effets fixes) : 14 \%. Les dummies aéroport capturent les
    spécificités locales (climat, géographie).
    \item \textbf{Distance} : 9 \%. \texttt{dist\_current}, \texttt{dist\_cum\_mean},
    \texttt{dist\_diff}.
    \item \textbf{Intensité} : 5 \%. \texttt{amp\_abs}, polarité, \% négatifs.
\end{itemize}

Les 4 coefficients individuels les plus forts :
\begin{itemize}
    \item \texttt{rolling\_gap\_5 = +0,354} : $\lambda \times e^{+0,354} \approx \times 1{,}42$.
    Un grand \texttt{rolling\_gap\_5} allonge le temps caractéristique de 42 \%.
    \item \texttt{cg\_rank = $-$0,373} : signe négatif. Plus on est tard dans l'orage (gros
    \texttt{cg\_rank}), plus $T_q$ est court — l'orage s'épuise.
    \item \texttt{rolling\_gap\_3 = +0,255} : signal plus court mais cohérent avec
    \texttt{rolling\_gap\_5}.
    \item \texttt{month\_cos} et \texttt{month\_sin} : saisonnalité, 5\textsuperscript{e} et
    6\textsuperscript{e} en importance.
\end{itemize}

Le message : c'est \textbf{le rythme passé} de l'orage qui est de loin le signal le plus
prédictif. La saison et l'aéroport contribuent modestement. La distance et l'intensité sont
secondaires. \fg{}
\end{dire}

\begin{compr}
\textbf{Pourquoi le rythme domine ?} Parce qu'un orage est un système physique avec de
l'inertie. Si les 5 derniers CG sont arrivés en rafale (\texttt{rolling\_gap\_5} petit), il y
a très probablement encore de l'activité électrique en cours dans le nuage — donc le prochain
CG arrive vite. Inversement, si \texttt{rolling\_gap\_5} est grand, l'orage est dans une phase
de pause / extinction.

\textbf{Comment ces 57 \% sont calculés} : on prend les coefficients $\beta$ standardisés
(c'est-à-dire après scaling des features), on somme leur valeur absolue par catégorie, et on
normalise pour avoir 100 \% au total. Méthode standard pour l'importance dans un modèle
linéaire.

\textbf{Le signe des coefficients} :
\begin{itemize}
    \item Positif (\texttt{rolling\_gap\_5 = +0,354}) : allonge $\lambda$ donc allonge $T_q$.
    \item Négatif (\texttt{cg\_rank = $-$0,373}) : raccourcit $\lambda$ donc raccourcit $T_q$.
\end{itemize}
\end{compr}

\begin{piege}
\textbf{Q : Pourquoi la distance ne fait que 9 \% ? Intuitivement, un orage qui s'éloigne devrait
être moins menaçant.}

R : Bonne question. La distance est utile mais elle est \emph{déjà partiellement capturée par le
rythme} : un orage qui s'éloigne fait moins d'éclairs en zone d'alerte, donc \texttt{cg\_rank}
plafonne et \texttt{rolling\_gap\_5} augmente. Donc l'effet \og distance \fg{} se retrouve en
grande partie dans les variables de rythme. La distance pure ajoute du signal marginal,
notamment quand l'orage est en bordure de la zone 20 km.

\textbf{Q : Pourquoi pas plus d'importance à la saisonnalité ? Les orages d'été sont différents.}

R : Ils SONT différents en \emph{nombre} (beaucoup plus d'orages en été), mais conditionnellement
à \emph{être dans une alerte}, leur dynamique intra-alerte n'est pas si différente. Le rythme
local domine.
\end{piege}

\begin{transit}
\og Le modèle est bon, mais avant de le déployer, il faut vérifier qu'il est bien calibré.
Voilà le piège, slide 10. \fg{}
\end{transit}

% ==========================================================================
\section{Slide 10 — Le piège : on annonce 1 \%, on a 6 \% en réalité}

\begin{dire}
\og On a un modèle qui prédit $T_q$. C'est bien, mais avant de le déployer, il faut vérifier
qu'il est \textbf{bien calibré}.

\textbf{Calibration} : si j'annonce 1 \% de risque, j'ai \emph{vraiment} 1 \% de risque, pas 6 \%.

Sur nos données test, le modèle n'est PAS calibré directement :
\begin{itemize}
    \item Quand on annonce 10 \% de risque, on en a 21 \% en réalité.
    \item Quand on annonce 5 \%, on en a 14 \%.
    \item Quand on annonce 1 \%, on en a 6 \%.
\end{itemize}
C'est entre 2 et 6 fois plus que ce qu'on annonce. Le modèle est \textbf{trop optimiste sur les
extrêmes}.

Pourquoi ? Parce que la queue de la loi de Weibull décroît très vite, en
$\exp(-(t/\lambda)^k)$. Mais les \emph{vrais} longs silences entre éclairs sont
\emph{encore plus longs} que ce que Weibull prédit. On est sous-calibré.

Donc on \textbf{ne peut pas déployer le modèle tel quel}. Il faut une étape de calibration
empirique post-hoc.

Mon binôme prend le relais pour expliquer comment on corrige ce piège. \fg{}
\end{dire}

\begin{compr}
\textbf{Calibration} : c'est la promesse que les probabilités annoncées correspondent à la
fréquence empirique. Exemple Météo France : si le bulletin annonce \og 30 \% de pluie demain \fg{},
on attend qu'il pleuve effectivement \emph{environ 30 \%} des jours où ce bulletin sort.

\textbf{Notre cas} : on annonce \og T\_q tel que la proba de futur éclair = 1 \% \fg{}, mais en
réalité 6 \% des observations dépassent ce T\_q. Donc le modèle est trop optimiste sur les
extrêmes (longs silences).

\textbf{Pourquoi Weibull sous-calibre} : sa queue décroît exponentiellement, ce qui est trop
rapide pour modéliser les vrais orages — où des silences longs peuvent survenir avec une
probabilité non négligeable. Les vrais orages ont une queue \og heavy-tailed \fg{} plus lourde
que ce que Weibull capture.

\textbf{Heureusement}, on peut corriger ça \emph{ex-post} par une procédure simple (winsorization
+ monotonicité) qui est l'objet du slide 11, donc le sujet de mon binôme.
\end{compr}

\begin{piege}
\textbf{Q : Pourquoi le Weibull est-il sous-calibré ?}

R : Parce que la queue Weibull décroît comme une exponentielle. Or les vrais longs silences
intra-orage suivent plutôt une décroissance plus lente (heavy tail). Le Weibull \og sous-estime
les valeurs extrêmes \fg{} — donc ses quantiles élevés ($T_{99}$) sont trop courts par rapport
à la réalité.

\textbf{Q : Pourquoi pas changer de loi (Gamma, log-normale) au lieu de calibrer post-hoc ?}

R : On a essayé. La log-normale donne des résultats similaires en C-index mais ses quantiles
extrêmes sont moins stables. La distribution généralisée de Pareto (heavy-tailed) pourrait
mieux modéliser la queue mais lifelines ne la propose pas en natif. La calibration empirique
post-hoc sur Weibull s'est révélée plus rapide à mettre en œuvre et donne des résultats
robustes.
\end{piege}

\begin{transit}
\og Mon binôme va maintenant expliquer comment on corrige ce piège, puis présenter les
résultats opérationnels. \fg{}
\end{transit}

% ==========================================================================
\section{Vue globale du code — le pipeline en 7 étapes}
\label{sec:code}

\textbf{Tout le pipeline est centralisé dans \texttt{meteorage\_model.py}} (202 lignes,
réutilisé par les notebooks et par Streamlit). Le code est volontairement modulaire :
chaque étape est une fonction pure, testable indépendamment.

\begin{tcolorbox}[colback=gray!5, colframe=gray!50!black, title=\textbf{Le pipeline en une ligne}]
\texttt{cg = load\_train(path)} $\to$ \texttt{add\_features(cg)} $\to$ \texttt{add\_target(cg)} $\to$
\texttt{build\_model\_matrix(cg, scaler, [], fit\_scaler=True)} $\to$ \texttt{fit\_weibull(fit\_df)} $\to$
\texttt{compute\_robust\_calibration(aft, fit\_df)} $\to$ \texttt{predict\_T\_q\_calibrated(aft, X, scaling, q)}
\end{tcolorbox}

\subsection{Étape 1 — \texttt{load\_train} : filtrage et identifiants uniques}

\begin{lstlisting}[caption={Chargement et nettoyage. Le filtre Pise 2016 vient d'une différence de système d'enregistrement.}]
def load_train(path: str) -> pd.DataFrame:
    df = pd.read_csv(path, parse_dates=['date'])
    alerts = df.dropna(subset=['airport_alert_id']).copy()
    # On rend l'ID unique en concatenant aeroport + numero d'alerte
    alerts['airport_alert_id'] = alerts['airport_alert_id'].astype(int).astype(str)
    alerts['airport_alert_id'] = alerts['airport'] + '_' + alerts['airport_alert_id']
    # Exclure Pise 2016 (systeme d'enregistrement different)
    alerts['year'] = alerts['date'].dt.year
    mask_pise_2016 = (alerts['airport'] == 'Pise') & (alerts['year'] == 2016)
    alerts = alerts.loc[~mask_pise_2016].copy()
    # Garder uniquement les CG (nuage-sol, icloud=False)
    cg = alerts.loc[~alerts['icloud']].copy()
    cg = cg.sort_values(['airport', 'airport_alert_id', 'date']).reset_index(drop=True)
    return cg
\end{lstlisting}

\textbf{Points-clés} :
\begin{itemize}[leftmargin=*]
    \item \texttt{dropna(subset=['airport\_alert\_id'])} : on ne garde que les éclairs dans une alerte (= zone 20 km).
    \item Concat \texttt{aeroport\_id} : sinon \og 530 \fg{} peut désigner 5 alertes différentes (une par aéroport).
    \item Exclusion Pise 2016 : justifiée techniquement (différent enregistrement IC).
    \item \texttt{icloud=False} : on ne garde que les CG, parce que le modèle prédit le temps jusqu'au prochain \emph{CG} (= ce qui ferme l'alerte).
\end{itemize}

\subsection{Étape 2 — \texttt{add\_features} : les 17 variables contextuelles}

\begin{lstlisting}[caption={Construction des features. Toutes utilisent \texttt{groupby + transform} pour rester intra-alerte.}]
def add_features(cg):
    cg = cg.copy()
    # 1) Saisonnalite cyclique (sin/cos)
    cg['month_sin'] = np.sin(2 * np.pi * cg['date'].dt.month / 12)
    cg['month_cos'] = np.cos(2 * np.pi * cg['date'].dt.month / 12)
    cg['hour_sin']  = np.sin(2 * np.pi * cg['date'].dt.hour / 24)
    cg['hour_cos']  = np.cos(2 * np.pi * cg['date'].dt.hour / 24)

    # 2) Rythme passe (intra-alerte avec groupby)
    grp = cg.groupby(['airport', 'airport_alert_id'], sort=False)
    cg['time_since_start'] = (cg['date'] - grp['date'].transform('first')).dt.total_seconds() / 60
    cg['prev_gap']         = grp['date'].diff().dt.total_seconds().div(60).fillna(0)
    cg['cg_rank']          = grp.cumcount() + 1
    cg['n_cg_alert']       = grp['cg_rank'].transform('max')
    cg['rolling_gap_3']    = grp['prev_gap'].transform(lambda x: x.rolling(3, min_periods=1).mean())
    cg['rolling_gap_5']    = grp['prev_gap'].transform(lambda x: x.rolling(5, min_periods=1).mean())
    cg['trend_gap']        = (cg['prev_gap'] / cg['rolling_gap_5'].replace(0, np.nan)).fillna(1)

    # 3) Distance / intensite (cumul intra-alerte avec expanding)
    cg['dist_current']  = cg['dist']
    cg['dist_cum_mean'] = grp['dist'].transform(lambda x: x.expanding().mean())
    cg['dist_diff']     = grp['dist'].diff().fillna(0)
    cg['amp_abs']       = cg['amplitude'].abs()
    cg['amp_cum_mean']  = grp['amp_abs'].transform(lambda x: x.expanding().mean())
    cg['is_negative']   = (cg['amplitude'] < 0).astype(int)
    cg['pct_neg_cum']   = grp['is_negative'].transform(lambda x: x.expanding().mean())
    return cg
\end{lstlisting}

\textbf{Pourquoi sin/cos pour mois/heure ?} L'encodage cyclique évite la fausse discontinuité
entre décembre (mois 12) et janvier (mois 1). Avec un encodage ordinaire, le modèle penserait
que ces deux mois sont \og loin \fg{}, alors qu'ils sont voisins dans le cycle annuel.

\textbf{Pourquoi \texttt{expanding} et pas \texttt{rolling} pour les cumuls de distance ?}
\texttt{expanding} accumule depuis le \emph{début de l'alerte}, ce qui capture la tendance
globale de l'orage. \texttt{rolling} fenêtre les $k$ derniers points, ce qui capture une
dynamique locale. On utilise les deux selon ce qu'on veut mesurer.

\textbf{Aucune fuite (\og leak \fg{})} : toutes les opérations sont causales — elles ne
regardent que les CG \emph{déjà arrivés}, jamais le futur. \texttt{groupby + transform}
garantit qu'on reste intra-alerte (pas de fuite inter-alerte).

\subsection{Étape 3 — \texttt{add\_target} : censure native à 30 min}

\begin{lstlisting}[caption={C'est ici qu'on encode la survie. duration et event sont les deux colonnes que lifelines attend.}]
def add_target(cg, is_last_col):
    cg = cg.copy()
    grp = cg.groupby(['airport', 'airport_alert_id'], sort=False)
    # Cible = temps jusqu'au prochain CG
    cg['target'] = grp['time_since_start'].shift(-1) - cg['time_since_start']
    cg['is_last'] = cg['cg_rank'] == cg['n_cg_alert']
    cg['duration'] = cg['target'].copy()
    cg['event'] = (~cg['is_last']).astype(int)
    # CENSURE : pour le dernier CG d'une alerte, on n'observe pas le prochain CG.
    # On sait juste qu'aucun CG n'est arrive pendant 30 min.
    cg.loc[cg['is_last'], 'duration'] = 30.0
    return cg
\end{lstlisting}

\textbf{La ligne magique} : \texttt{cg.loc[cg['is\_last'], 'duration'] = 30.0}. Pour le dernier
CG d'une alerte, on ne sait pas combien de temps il a fallu avant le \emph{vrai} prochain CG
(il y en a un, plus tard, mais ça démarre une nouvelle alerte). On sait juste qu'il n'y a rien
eu pendant au moins 30 min. C'est exactement une observation \textbf{censurée à droite} :
\texttt{event=0}, \texttt{duration=30}. lifelines intègre ça nativement dans la vraisemblance.

\subsection{Étape 4 — \texttt{build\_model\_matrix} : standardisation + dummies aéroport}

\begin{lstlisting}[caption={Préparation finale. fit\_scaler=True uniquement sur train.}]
def build_model_matrix(cg, scaler, apt_cols, fit_scaler=False):
    # On filtre les alertes avec moins de 3 CG (pas assez de contexte) et duration<=0
    model_df = cg[cg['n_cg_alert'] >= 3].copy()
    model_df = model_df[model_df['duration'] > 0].copy()
    for col in FEATURES:
        model_df[col] = model_df[col].replace([np.inf, -np.inf], np.nan).fillna(0)

    # 1) Standardisation des 17 features (moyenne 0, ecart-type 1)
    if fit_scaler:
        scaler.fit(model_df[FEATURES])
    model_df_sc = model_df.copy()
    model_df_sc[FEATURES] = scaler.transform(model_df[FEATURES])

    # 2) Dummies aeroport (one-hot, drop_first pour eviter colinearite)
    apt_dum = pd.get_dummies(model_df_sc['airport'], prefix='apt', drop_first=True).astype(float)
    if fit_scaler:
        apt_cols = list(apt_dum.columns)
    else:
        # En eval : on aligne les colonnes pour avoir le meme set qu'au train
        for col in apt_cols:
            if col not in apt_dum.columns:
                apt_dum[col] = 0.0
        apt_dum = apt_dum[apt_cols]

    fit_df = pd.concat([
        model_df_sc[FEATURES + ['duration', 'event']].reset_index(drop=True),
        apt_dum.reset_index(drop=True),
    ], axis=1)
    fit_df = fit_df.replace([np.inf, -np.inf], np.nan).fillna(0)
    fit_df = fit_df[fit_df['duration'] > 0]
    return model_df_sc.reset_index(drop=True), fit_df, apt_cols
\end{lstlisting}

\textbf{Standardisation} : indispensable pour qu'aucune feature ne domine artificiellement
(ex.\ \texttt{time\_since\_start} en minutes vs \texttt{is\_negative} qui vaut 0 ou 1).
\textbf{Le scaler est fit \emph{uniquement} sur train} (paramètre \texttt{fit\_scaler=True})
et utilisé en \og transform-only \fg{} sur eval — c'est la règle d'or pour éviter la fuite.

\textbf{Dummies aéroport, \texttt{drop\_first=True}} : on crée 4 colonnes binaires pour les 5
aéroports (Bastia, Biarritz, Nantes, Pise vs référence Ajaccio). \texttt{drop\_first} évite
la colinéarité parfaite (sinon les 5 colonnes somment à 1 = $\beta_0$, le modèle ne peut pas
distinguer).

\subsection{Étape 5 — \texttt{fit\_weibull} : 3 lignes pour le MLE avec censure}

\begin{lstlisting}[caption={lifelines fait tout le travail.}]
def fit_weibull(train_fit, penalizer=0.05):
    aft = WeibullAFTFitter(penalizer=penalizer, l1_ratio=0.0)
    aft.fit(train_fit, duration_col='duration', event_col='event')
    return aft
\end{lstlisting}

\textbf{Ce que ces 3 lignes font sous le capot :}
\begin{itemize}[leftmargin=*]
    \item Maximisation numérique (Newton-Raphson) de la log-vraisemblance avec censure.
    \item Régularisation L2 ridge avec \texttt{penalizer=0.05}. Évite l'overfit, plus stable
    numériquement avec 17+ features.
    \item \texttt{l1\_ratio=0.0} : 100 \% ridge, pas de sparsity Lasso (on garde tous les
    coefficients).
\end{itemize}

\textbf{Sortie} : un objet \texttt{aft} qui expose \texttt{predict\_percentile(X, p=1-q)} pour
calculer $T_q$, \texttt{predict\_survival\_function(X, times=[30])} pour $S(30|X)$, etc.

\subsection{Étape 6 — \texttt{compute\_robust\_calibration} : winsorize + monotonicité}

\begin{lstlisting}[caption={Calibration ex-post pour aligner risque annonce et risque reel.}]
def compute_robust_calibration(aft, fit_df, q_levels=[0.90, 0.95, 0.99]):
    evt = fit_df[fit_df['event'] == 1].copy()
    y_real = evt['duration'].values

    c_q_winsorized = {}
    for q in q_levels:
        # Tq predit par Weibull pour le quantile q
        Tq = aft.predict_percentile(evt, p=1.0 - q).values
        # Residu = ratio entre temps reel et temps predit
        residuals = y_real / np.where(Tq > 1e-9, Tq, 1e-9)
        # Winsorize au P95 pour neutraliser les outliers
        p95 = np.quantile(residuals, 0.95)
        residuals_wins = np.minimum(residuals, p95)
        # Coefficient c_q = quantile q des residus winsorises
        c_q_winsorized[q] = float(np.quantile(residuals_wins, q))

    # Imposer la monotonicite : c_99 >= c_95 >= c_90
    scaling = {}
    prev_c = 0.0
    for q in sorted(q_levels):
        c_q_mono = max(c_q_winsorized[q], prev_c)
        scaling[q] = c_q_mono
        prev_c = c_q_mono
    return scaling
\end{lstlisting}

\subsection{Étape 7 — Prédiction calibrée et confidence}

\begin{lstlisting}[caption={Utilisation en production. Cap a 30 min = filet de securite operationnel.}]
def predict_T_q_calibrated(aft, X, scaling, q):
    Tq_raw = aft.predict_percentile(X, p=1.0 - q).values
    return scaling[q] * Tq_raw

# Pour la confidence (utilisee dans le protocole officiel) :
# confidence = S(30 | X) = exp(-(30 / lambda(X))^k)
survival_df = aft.predict_survival_function(X, times=[30.0])
confidence = survival_df.iloc[0].values

# Cap operationnel a 30 min : on ne veut jamais lever PLUS TARD que la baseline
Tq_cap = np.clip(scaling[q] * Tq_raw, 0, 30)
\end{lstlisting}

% ==========================================================================
\section{Q/R techniques anticipées (sur le code et le modèle)}
\label{sec:qa-tech}

\subsection*{Sur la bibliothèque et l'algorithme}

\begin{piege}
\textbf{Q : Pourquoi \texttt{lifelines} et pas une implémentation maison ?}

R : \texttt{lifelines} est la référence Python pour l'analyse de survie (Cameron Davidson-Pilon).
Implémentation maison : on aurait dû écrire la log-vraisemblance avec censure, le gradient,
gérer les contraintes ($\lambda > 0, k > 0$), et l'optimisation Newton-Raphson. lifelines
fait tout ça avec une API propre et testée.
\end{piege}

\begin{piege}
\textbf{Q : Comment lifelines gère la censure dans la vraisemblance ?}

R : Pour une observation $(t_i, \delta_i, X_i)$ avec $\delta_i \in \{0, 1\}$ :
\begin{itemize}
    \item Si $\delta_i = 1$ (événement observé) : contribution = $\log f(t_i \mid X_i)$ (densité)
    \item Si $\delta_i = 0$ (censure) : contribution = $\log S(t_i \mid X_i)$ (survie)
\end{itemize}
La log-vraisemblance totale est la somme. lifelines optimise par Newton-Raphson sur les
paramètres $(\beta, \log k)$ — le $\log$ assure que $k > 0$ implicitement.
\end{piege}

\begin{piege}
\textbf{Q : Pourquoi \texttt{penalizer=0.05} et pas 0 ou 0.1 ?}

R : Choix empirique. À 0 (pas de pénalisation), le modèle est plus instable numériquement avec
17 features + 4 dummies. À 0.1, il est trop régularisé et perd un peu de pouvoir prédictif
(C-index baisse). 0.05 est un bon compromis trouvé par cross-validation rapide. C'est une
pénalité ridge $\lambda \sum \beta_j^2$ ajoutée à la log-vraisemblance.
\end{piege}

\begin{piege}
\textbf{Q : Pourquoi pas un gradient boosting (XGBoost / LightGBM) ?}

R : Trois raisons :
\begin{enumerate}
    \item \textbf{La censure} : les GBM classiques ne savent pas gérer nativement les données
    censurées (un éclair où on sait juste qu'il n'y a rien eu pendant 30 min). Weibull AFT oui.
    \item \textbf{Les quantiles analytiques} : on a la formule $T_q = \lambda \cdot (-\ln q)^{1/k}$.
    Un GBM nous donnerait juste une médiane.
    \item \textbf{Interprétabilité} : on lit directement l'effet de chaque variable via $\beta_j$.
    Un GBM est plus boîte noire.
\end{enumerate}
On a quand même testé un GBM en hybride (cellule 16 du notebook \texttt{weibull\_final}) à
titre diagnostic — Weibull AFT calibré reste meilleur sur notre C-index.
\end{piege}

\begin{piege}
\textbf{Q : Comment \texttt{predict\_percentile(p)} fonctionne ?}

R : Pour un Weibull paramétré par $\lambda(X)$ et $k$ :
\[
S(t \mid X) = \exp(-(t/\lambda)^k) = p \quad \Rightarrow \quad t = \lambda(X) \cdot (-\ln p)^{1/k}
\]
\texttt{predict\_percentile(X, p)} retourne ce $t$. Donc pour avoir le quantile $q$ (qui couvre
$q$ \% des cas), on appelle avec \texttt{p = 1 - q}. C'est pour ça qu'on a $p = 1 - q$ partout.
\end{piege}

\subsection*{Sur les features et la cible}

\begin{piege}
\textbf{Q : Pourquoi exclure les alertes avec \texttt{n\_cg\_alert < 3} ?}

R : Pour les features de rythme (rolling\_gap\_3, rolling\_gap\_5, trend\_gap), il faut au
moins 3 points pour avoir un signal exploitable. Avec 1 ou 2 CG, les rolling sont dégénérés
(NaN ou identiques à \texttt{prev\_gap}). On préfère exclure ces cas en entraînement —
ils représentent moins de 5 \% des alertes et ajoutent du bruit.
\end{piege}

\begin{piege}
\textbf{Q : Pourquoi pas de \texttt{prev\_gap} pour le tout premier CG ?}

R : Effectivement le premier CG n'a pas de prédécesseur. On met \texttt{prev\_gap = 0} et
\texttt{cg\_rank = 1}. Ces cas sont rares (1/N par alerte) et leur contribution au fit est
faible. Une alternative serait d'utiliser le gap inter-alerte (temps depuis la fin de l'alerte
précédente), mais on ne l'a pas exploité ici.
\end{piege}

\begin{piege}
\textbf{Q : Comment gérer \texttt{duration = 0} ? (deux éclairs simultanés)}

R : On les exclut avec \texttt{model\_df[model\_df['duration'] > 0]} dans
\texttt{build\_model\_matrix}. Causes possibles : précision temporelle de Météorage à la
seconde, deux éclairs détectés au même horodatage. Les exclure évite la singularité dans le
log-vraisemblance ($\log(0) = -\infty$).
\end{piege}

\begin{piege}
\textbf{Q : Pourquoi \texttt{drop\_first=True} pour les dummies aéroport ?}

R : Pour éviter la colinéarité parfaite. Avec 5 aéroports, créer 5 dummies (\texttt{apt\_Ajaccio},
\texttt{apt\_Bastia}, ...) donne 5 colonnes qui somment toujours à 1. Combiné à
l'intercept $\beta_0$, ça donne une matrice de design singulière. \texttt{drop\_first} retire
une colonne (Ajaccio devient la \og référence \fg{}), les coefficients $\beta$ des 4 autres
mesurent l'écart à Ajaccio.
\end{piege}

\subsection*{Sur le scaler et l'eval}

\begin{piege}
\textbf{Q : Comment le scaler est-il appliqué en eval ?}

R : On fit le scaler une fois sur train (\texttt{fit\_scaler=True}), il apprend la moyenne et
l'écart-type de chaque feature \emph{sur train}. En eval, on l'utilise en mode transform
uniquement (\texttt{fit\_scaler=False}) — il applique les mêmes paramètres train aux données
eval. C'est la règle d'or pour éviter la fuite : \textbf{aucune statistique du test ne doit
fuiter dans le fit}.
\end{piege}

\begin{piege}
\textbf{Q : Et si un aéroport apparaît en eval mais pas en train ?}

R : Géré par le code \texttt{build\_model\_matrix} :
\begin{lstlisting}[basicstyle=\ttfamily\scriptsize]
for col in apt_cols:
    if col not in apt_dum.columns:
        apt_dum[col] = 0.0
apt_dum = apt_dum[apt_cols]
\end{lstlisting}
Si un aéroport nouveau apparaît, sa dummy n'existe pas dans \texttt{apt\_cols} (qui vient
du train), donc il devient implicitement la \og référence \fg{} (Ajaccio dans notre cas).
C'est imparfait mais ça ne crash pas. Sur eval 2023-25, les 5 aéroports sont les mêmes
qu'en train, donc ce cas n'arrive pas.
\end{piege}

\begin{piege}
\textbf{Q : Le scaler influence-t-il la calibration ?}

R : Non. La calibration utilise $T_q$ prédit par le modèle (déjà sur l'échelle des minutes,
pas standardisée) et les vraies durées en minutes. Le scaler intervient \emph{avant} le fit
(sur les features X), pas sur la cible \texttt{duration}. La calibration travaille uniquement
sur des durées en minutes, donc le scaler est transparent.
\end{piege}

\subsection*{Sur la calibration}

\begin{piege}
\textbf{Q : Pourquoi winsoriser au P95 et pas un autre seuil ?}

R : Choix empirique. P95 retire les ~5 \% d'outliers extrêmes (en pratique des observations
où l'orage a dévié totalement de ce que prédit Weibull, souvent à cause d'une transition de
système orageux). Un seuil plus bas (P90) corrigerait trop fort. Un seuil plus haut (P99)
laisserait passer les outliers. P95 est le compromis standard en stats robustes.
\end{piege}

\begin{piege}
\textbf{Q : Pourquoi $c_{99}^{\text{wins}} < c_{95}$ après winsorization ?}

R : Effet artificiel de la winsorization. Quand on tronque la queue à P95, les valeurs au-dessus
du P95 deviennent toutes égales. Le quantile $q=0.99$ \emph{après} winsorization est donc
le P95 \emph{avant} winsorization. Si Weibull surestime déjà cette partie de la queue, $c_{99}$
peut devenir plus petit que $c_{95}$. L'imposition de monotonicité ($c_q = \max(c_q, c_{q-1})$)
corrige cet artefact.
\end{piege}

\begin{piege}
\textbf{Q : Pourquoi la calibration est-elle faite sur le test et pas sur un fold dédié ?}

R : Limite assumée. L'idéal est \textbf{3 splits} : train pour fit, calibration pour estimer $c_q$,
test pour évaluation finale. On a fait 2 splits par contrainte de volume (54k CG). En
production, on referait avec 3 splits. La winsorization protège partiellement de l'overfit
à la calibration en neutralisant les outliers.
\end{piege}

\subsection*{Sur le critère de risque et l'évaluation}

\begin{piege}
\textbf{Q : Comment le C-index est-il calculé ?}

R : Pour toute paire $(i, j)$ d'observations non-censurées dans le test, on regarde si
le modèle prédit la même hiérarchie que la réalité. \texttt{concordance\_index(durations,
predictions, events)} de lifelines fait ça en $O(N^2)$ ou avec un index spécial en $O(N \log N)$.
Notre prédiction utilisée : \texttt{aft.predict\_median(test\_df)} (la médiane prédite).
\end{piege}

\begin{piege}
\textbf{Q : Pourquoi cap à 30 min ? Vous bridez le modèle.}

R : C'est un \og safety net \fg{}. La règle Météorage 30 min est calibrée pour 1 \% de risque
sur 20 km. Si notre modèle propose $T_q > 30$ min, on préfère garder 30 min — on ne veut
\emph{jamais être moins prudent} que la baseline. Décision opérationnelle, pas mathématique.
Sans le cap, $T_{99}$ peut dépasser 90 min sur des cas extrêmes.
\end{piege}

\begin{piege}
\textbf{Q : La reproductibilité ? Random seed ?}

R : Le fit Weibull est déterministe (optimisation convexe avec pénalité ridge, point initial
fixé par lifelines). Pas besoin de seed. \texttt{predict\_percentile} est aussi déterministe.
Donc en relançant le notebook on obtient exactement les mêmes chiffres — c'est testé.
\end{piege}

\subsection*{Sur le déploiement et la production}

\begin{piege}
\textbf{Q : Combien de temps prend une prédiction ?}

R : Une prédiction $T_q(X)$ = un produit matriciel sur un vecteur de 17+4 dimensions =
quelques microsecondes. Le coût opérationnel est \emph{négligeable} — l'inférence est
1000$\times$ plus rapide qu'un appel API externe. Le bottleneck en production serait
l'agrégation Météorage (qui détecte les éclairs en temps réel).
\end{piege}

\begin{piege}
\textbf{Q : Comment vous déploieriez en production ?}

R :
\begin{enumerate}
    \item Sérialisation du modèle avec \texttt{joblib.dump(aft, scaler, scaling)} (en
    pratique \texttt{meteorage\_artifact.pkl} déjà présent dans le repo).
    \item Service Python (FastAPI ou Flask) qui reçoit le contexte courant et retourne $T_q$.
    \item Side-by-side avec la baseline Météorage pendant 3-6 mois (\og shadow mode \fg{}).
    \item Mesure A/B des gains réels en heures d'opération sauvées.
    \item Si OK, bascule progressive aéroport par aéroport.
    \item Re-fit annuel automatisé sur les nouvelles données.
\end{enumerate}
\end{piege}

\begin{piege}
\textbf{Q : Et si Météorage change ses détecteurs ?}

R : Le modèle dépend des caractéristiques de détection (sensibilité, fréquence d'IC vs CG).
Un changement de capteurs créerait une dérive — il faudrait re-fit. Pour le détecter,
on monitorerait la distribution des features (data drift) avec des tests KS ou PSI.
\end{piege}

\begin{piege}
\textbf{Q : Le modèle est-il transférable à d'autres aéroports ?}

R : Partiellement. Les coefficients $\beta$ sur les 17 features (rythme, distance, intensité,
saison) sont \emph{universels} — un orage actif dans la Bresse se comporte comme un orage
actif à Ajaccio sur ces variables. Les dummies aéroport capturent les spécificités locales
(climat, géographie). Pour un nouvel aéroport, on aurait besoin de 6-12 mois de données pour
estimer son effet local. Ou bien : utiliser le modèle sans dummy aéroport (= généraliste) et
accepter une perte de C-index de quelques points.
\end{piege}

\begin{piege}
\textbf{Q : Avez-vous testé d'autres distributions (Gamma, log-normale, exponentielle) ?}

R : Brièvement. L'exponentielle est un cas particulier de Weibull avec $k = 1$ — on l'a
rejetée par AIC. Le log-normale donne des résultats similaires en C-index mais ses quantiles
extrêmes sont moins stables. La distribution généralisée de Pareto (queue heavy-tailed)
pourrait améliorer la calibration mais lifelines ne la propose pas en natif. Weibull est
le standard et permet une calibration ex-post efficace.
\end{piege}

% ==========================================================================
\section*{Mémo express — 1 minute avant l'oral}

\begin{tcolorbox}[colback=red!4, colframe=red!50!black, title=\textbf{Les 5 phrases-clés (slides 1-10)}]
\begin{enumerate}
    \item \textbf{Problème} : la règle 30 min Météorage est aveugle au contexte (slide 1).
    \item \textbf{Approche} : Weibull AFT prédit $T_q$ adapté au contexte par éclair (slides 2-3).
    \item \textbf{Modèle} : Weibull AFT ($k = 0{,}697$ → hazard décroissant) entraîné par MLE
    avec gestion native de la censure 30 min (slides 4-5).
    \item \textbf{17 features} dont le \textbf{rythme} concentre 57 \% du pouvoir prédictif ;
    C-index = 0,74 vs 0,55 pour les baselines simples (slides 6-9).
    \item \textbf{Limite} : modèle sous-calibré ($1\%$ annoncé → $6\%$ réel) → mon binôme prend
    le relais sur la calibration et les résultats (slide 10).
\end{enumerate}
\end{tcolorbox}

\begin{tcolorbox}[colback=blue!4, colframe=blue!50!black, title=\textbf{Lancer la démo}]
\texttt{streamlit run app.py}

(à lancer dans un terminal positionné dans \texttt{oral\_projet/})
\end{tcolorbox}

\end{document}
